Nach dem Bremsen ist die Geschwindigkeit
\begin{align}
 v &= \vVeF = \vVa+\left(\aVb\right)\cdot\tVb = \vVe \approx \result{\vVeP{0}{2}}
\end{align}
Die Geschwindigkeit nimmt gleichmässig ab, deshalb ist die Durchschnittsgeschwindigkeit der Mittelwert aus Anfangs- und Endgeschwindigkeit:
\begin{align}
 \bar v &= \vVmF = \frac{\vVa+\vVe}{2} = \vVm \approx \result{\vVmP{0}{2}}
\end{align}
Der Bremsweg ist
\begin{align}
 s &= \sVbF = \vVm\cdot\tVb = \sVb \approx \result{\sVbP{0}{2}}
\end{align}
Mit der allgemeinen Formel für die gleichmässig beschleunigte Bewegung erhält man dasselbe:
\begin{align}
 s &= \sVaF \\
   &= \vVa\cdot\tVb+\frac{1}{2}\cdot\left(\aVb\right)\cdot\left(\tVb\right)^2 = \sVa \approx \result{\sVaP{0}{2}}
\end{align}
